\section{La pila de problemas de eficiencia energética: más allá de los FLOPS hay tres muros}

Al bajar de la difusión de las aplicaciones a la capa inferior, lo primero que aparece es la eficiencia energética. El ponente enumera cuatro frentes: semiconductor, memory, interconnect y cooling. Lo que en realidad quiere decir es que el throughput de un AI system está limitado por el recurso más lento y que los FLOPS pico son solo una cota superior. Esta sección unifica estas restricciones con un razonamiento de tipo roofline, para no llamar «faltan GPU» a todos los cuellos de botella.

\subsection{Von Neumann bottleneck y data movement}

El \term{von Neumann bottleneck} (cuello de botella de von Neumann) se refiere a que las unidades de ejecución del procesador están separadas del almacenamiento, por lo que los datos y las instrucciones deben moverse entre la memory hierarchy y el compute. En muchos AI kernels, la energía y el tiempo no se consumen en las multiplicaciones y sumas en sí, sino en leer weights, KV cache y activations, y en el intercambio entre dispositivos. Para una workload puede escribirse una cota inferior simplificada
\[
T_{step}\geq \max\left(\frac{F}{C},\frac{B_m}{BW_m},\frac{B_n}{BW_n}\right),
\]
donde $F$ es la cantidad de operaciones requerida y $C$ es el throughput de cómputo efectivo; $B_m$ y $BW_m$ son, respectivamente, el memory traffic y el memory bandwidth; $B_n$ y $BW_n$ son el network traffic y el interconnect bandwidth. La fórmula indica que, si solo se aumenta $C$, los otros dos términos todavía pueden dominar el tiempo.

\lecturefigure{02-efficiency-stack.png}{La eficiencia energética de la IA exige revisar a la vez semiconductor, memory, interconnect y facility.}{Subtítulos locales 06:20--12:00; redibujo conceptual.}

\begin{knowledgebox}{Lectura de la figura: las cuatro capas no son cuatro proyectos independientes}
Un nuevo semiconductor puede aumentar la rack density y, con ello, modificar el cooling y el suministro eléctrico; una \term{High Bandwidth Memory} (HBM, DRAM apilada de alto ancho de banda) más grande puede reducir la comunicación entre tarjetas, pero eleva el costo de encapsulado; el optical interconnect puede disminuir las pérdidas en la transmisión a larga distancia, pero introduce nuevos dispositivos y una nueva pila de software. La pregunta correcta es cuántos SLO-compliant tokens per watt obtiene la workload de extremo a extremo, no si el indicador pico de algún dispositivo luce mejor.
\end{knowledgebox}

\subsection{Qué resuelve y qué no resuelve el compound semiconductor}

Un \term{compound semiconductor} (semiconductor compuesto) está formado por dos o más elementos, por ejemplo GaN (gallium nitride, nitruro de galio) o SiC (silicon carbide, carburo de silicio). En alto voltaje, alta frecuencia y conversión de potencia pueden ofrecer mejor campo de ruptura o mejor eficiencia, por lo que son adecuados para fuentes de alimentación, radiofrecuencia y dispositivos optoelectrónicos; sin embargo, «migrar por completo desde el silicon» no es un interruptor único. Los chips lógicos, la memory, la power electronics y la photonics tienen requisitos distintos de materiales, rendimiento de fabricación y costo.

\begin{warningbox}{La analogía de dispositivos de la clase no debe tomarse como hoja de ruta de producto}
Los subtítulos mezclan switching voltage, leakage y compound materials en una sola respuesta rápida. Estas notas conservan la afirmación de que «la capa de dispositivos todavía tiene margen de eficiencia energética», pero no presentan ningún material como sustituto general de la GPU. La viabilidad técnica depende además de process integration, yield, packaging, toolchain y economías de escala.
\end{warningbox}

\subsection{Resumen de la sección}

La eficiencia energética de la infraestructura de IA es una propiedad del sistema de extremo a extremo. Si cualquiera de las capas (aritmética, memory, network o facility) se queda corta, un accelerator costoso puede quedar en espera; tanto la investigación como las compras deberían usar como indicador común el useful output de workloads reales.
